\documentclass[11pt,a4paper]{article}
\usepackage[a4paper,top=10mm,bottom=10mm,left=14mm,right=14mm]{geometry}
\usepackage[T1]{fontenc}
\usepackage{amsmath}
\usepackage{amssymb}
\usepackage{enumitem}
\usepackage{multicol}
\usepackage{array}
\usepackage{xcolor}
\usepackage{tikz}
\usepackage[most]{tcolorbox}
\usetikzlibrary{arrows.meta,decorations.pathreplacing}

% ================================================================
% Picture sequences and the nth term: an extension sheet.
%
% A follow on from timesTablesToNthTerms, for pupils who have met
% times table sequences such as 3n. Part 1 finds the nth term of a
% linear picture sequence by splitting each pattern into the part
% that is the same every time and the groups that are added. A worked
% example is followed by the same layout with gaps, and then by
% questions with the layout taken away.
%
% Part 2 builds square numbers from squares of dots, and sequences
% made from a square with dots added or taken away. Part 3 has pixel
% art sequences to investigate: a rocket that grows, a heart that is
% enlarged, a pyramid that rearranges into a square and a frame around
% a square picture. It ends with pupils designing their own. The last
% side is a set of always, sometimes or never statements.
%
% Every picture is drawn from n by the macros below, so that the
% number of objects in each pattern is the value of the nth term.
% ================================================================

\pagestyle{empty}
\setlength{\parindent}{0pt}
\setlength{\parskip}{3pt}
\setlength{\multicolsep}{3pt}
\setlength{\columnsep}{14pt}
\setlist[enumerate,1]{label=\textbf{Q\arabic*.}, leftmargin=*, itemsep=4pt, topsep=3pt, parsep=2pt}
\setlist[enumerate,2]{label=(\alph*), leftmargin=*, itemsep=5pt, topsep=3pt, parsep=1pt}
\newcolumntype{C}[1]{>{\centering\arraybackslash}m{#1}}

% A gap for one number in a sentence.
\newcommand{\gap}{\underline{\hspace{0.95cm}}}
\newcommand{\ansline}[1][2.2cm]{\underline{\hspace{#1}}}
% The strut that sets the height of a table row to be written in.
\newcommand{\rowstrut}{\rule[-1.9mm]{0pt}{5.8mm}}

% The idea each part is built on, stated once and boxed, and the
% worked example. The boxes are tinted only faintly, and the title bar
% is pale with coloured text, so that the sheet is cheap to print.
\newtcolorbox{keyidea}{colback=orange!2,colframe=orange!60!black,boxrule=0.6pt,
  left=6pt,right=6pt,top=2pt,bottom=2pt}
\newtcolorbox{wex}[1]{colback=blue!2,colframe=blue!30,boxrule=0.6pt,
  colbacktitle=blue!8,coltitle=blue!50!black,
  left=6pt,right=6pt,top=3pt,bottom=4pt,
  title={\bfseries\small Worked example: #1},fonttitle=\small,
  toptitle=1pt,bottomtitle=1pt}

% The section headings, kept tight so the parts do not eat the page.
\newcommand{\parttitle}[2]{%
  \vspace{1pt}
  {\color{blue!45!black}\rule{\linewidth}{1pt}}\\[1pt]
  {\large\bfseries Part #1: #2}\\[-6pt]
  {\color{blue!45!black}\rule{\linewidth}{0.4pt}}
  \par\vspace{-5pt}}

% ----------------------------------------------------------------
% The table for a picture sequence, with one \picrow for each row.
% The last column holds the nth term. A \comparerow is a further row
% to compare the objects with, and a \filledrow is a row with its
% values written in, for the worked example.
\newenvironment{pictable}{%
  \setlength{\tabcolsep}{2pt}%
  \begin{tabular}{|C{1.85cm}|*{6}{C{0.74cm}|}C{1.75cm}|}\hline
  {\small Pattern $n$} & 1 & 2 & 3 & 4 & 5 & 10 & {\small $n$th term}\\\hline}
  {\end{tabular}}
\newcommand{\picrow}[1]{{\small #1} & & & & & & & \rowstrut\\\hline}
\newcommand{\comparerow}[1]{{\small #1} & & & & & & & \multicolumn{1}{c}{}\rowstrut\\\cline{1-7}}
\newcommand{\filledrow}[8]{{\small #1} & #2 & #3 & #4 & #5 & #6 & #7 & #8\rowstrut\\\hline}

% A picture sequence on the left and its table on the right.
\newcommand{\pictureandtable}[2]{%
  \begin{minipage}[c]{0.42\linewidth}#1\end{minipage}\hfill
  \begin{minipage}[c]{0.57\linewidth}\raggedleft #2\end{minipage}}

% The layout of the worked example, used again for the questions that
% follow it: \splitlayout{<patterns>}{<label>}{<pattern 3 split up>}
% {<table>}{<sentence>}
\newcommand{\splitlayout}[5]{%
  \begin{minipage}[c]{0.5\linewidth}#1\end{minipage}\hfill
  \begin{minipage}[c]{0.48\linewidth}\raggedleft{\small #2}\quad #3\end{minipage}\par\smallskip
  {\centering #4\par}\smallskip
  #5}

% ----------------------------------------------------------------
% Picture sequences. In each of these the argument is the pattern
% number n.
\tikzset{
  stick/.style={line width=1.1pt,cap=round,shorten >=1.3pt,shorten <=1.3pt,
    -{Circle[length=3.2pt,width=3.2pt]}},
  brace/.style={decorate,decoration={brace,mirror,amplitude=3pt},thin},
  lab/.style={font=\scriptsize,anchor=north,inner sep=1pt}}
\newcommand{\pic}[1]{\begin{tikzpicture}[baseline=(current bounding box.south)]#1\end{tikzpicture}}
% One matchstick from one point to another.
\newcommand{\stick}[2]{\draw[stick] (#1) -- (#2);}

% The 3 matchsticks that square k adds to a row of squares.
\newcommand{\squaregroup}[1]{%
  \stick{0.6*#1-0.6,0}{0.6*#1,0}\stick{0.6*#1-0.6,0.6}{0.6*#1,0.6}
  \stick{0.6*#1,0}{0.6*#1,0.6}}
% A row of n squares of matchsticks that share their sides.
\newcommand{\chainsquares}[1]{\pic{%
  \stick{0,0}{0,0.6}
  \foreach \k in {1,...,#1} {\squaregroup{\k}}}}
% The same row with the first matchstick and each group of 3 pulled apart.
\newcommand{\explodedsquares}[1]{\pic{%
  \stick{0,0}{0,0.6}
  \node[lab] at (0,-0.22) {1};
  \foreach \k in {1,...,#1} {
    \begin{scope}[shift={(0.3*\k,0)}]
      \squaregroup{\k}
      \draw[brace] (0.6*\k-0.6,-0.08) -- (0.6*\k,-0.08);
      \node[lab] at (0.6*\k-0.3,-0.22) {3};
    \end{scope}}}}

% The 5 matchsticks that house k adds to a row of houses.
\newcommand{\housegroup}[1]{%
  \stick{0.55*#1-0.55,0}{0.55*#1,0}\stick{0.55*#1-0.55,0.55}{0.55*#1,0.55}
  \stick{0.55*#1,0}{0.55*#1,0.55}
  \stick{0.55*#1-0.55,0.55}{0.55*#1-0.275,0.87}\stick{0.55*#1-0.275,0.87}{0.55*#1,0.55}}
% A row of n houses of matchsticks that share their walls.
\newcommand{\chainhouses}[1]{\pic{%
  \stick{0,0}{0,0.55}
  \foreach \k in {1,...,#1} {\housegroup{\k}}}}
% The same row with the first matchstick and each group pulled apart.
% The braces are left without numbers for pupils to write in.
\newcommand{\explodedhouses}[1]{\pic{%
  \stick{0,0}{0,0.55}
  \node[lab] at (0,-0.22) {\phantom{1}};
  \foreach \k in {1,...,#1} {
    \begin{scope}[shift={(0.3*\k,0)}]
      \housegroup{\k}
      \draw[brace] (0.55*\k-0.55,-0.08) -- (0.55*\k,-0.08);
      \node[lab] at (0.55*\k-0.275,-0.22) {\phantom{5}};
    \end{scope}}}}

% A row of n square tables pushed together. Each dot is a chair.
\newcommand{\tablesrow}[1]{\pic{%
  \draw[thick,sharp corners] (0,0) rectangle (0.5*#1,0.5);
  \foreach \k in {1,...,#1} {
    \ifnum\k<#1 \draw[thick] (0.5*\k,0) -- (0.5*\k,0.5);\fi
    \fill (0.5*\k-0.25,-0.17) circle (0.08);
    \fill (0.5*\k-0.25,0.67) circle (0.08);}
  \fill (-0.17,0.25) circle (0.08);
  \fill (0.5*#1+0.17,0.25) circle (0.08);}}

% A cross of dots: one centre dot and four arms of n dots.
\newcommand{\dotcross}[1]{\pic{%
  \fill (0,0) circle (0.065);
  \foreach \i in {1,...,#1} \foreach \a in {0,90,180,270}
    \fill (\a:0.22*\i) circle (0.065);}}

% An L of dots with n dots on each arm, sharing the corner dot.
\newcommand{\dotL}[1]{\pic{%
  \foreach \i in {1,...,#1} {\fill (0.25,0.25*\i) circle (0.065);}
  \ifnum#1>1 \foreach \i in {2,...,#1} {\fill (0.25*\i,0.25) circle (0.065);}\fi}}

% A square of n by n dots.
\newcommand{\dotsquare}[1]{\pic{%
  \foreach \i in {1,...,#1} \foreach \j in {1,...,#1}
    \fill (0.25*\i,0.25*\j) circle (0.06);}}
% A rectangle of #1 columns and #2 rows of dots.
\newcommand{\dotrect}[2]{\pic{%
  \foreach \i in {1,...,#1} \foreach \j in {1,...,#2}
    \fill (0.25*\i,0.25*\j) circle (0.06);}}
% A square of n by n dots with two more dots beside it.
\newcommand{\squareplustwo}[1]{\pic{%
  \foreach \i in {1,...,#1} \foreach \j in {1,...,#1}
    \fill (0.25*\i,0.25*\j) circle (0.06);
  \fill (0.25*#1+0.5,0.25) circle (0.06);
  \fill (0.25*#1+0.5,0.5) circle (0.06);}}
% Two squares of n by n dots.
\newcommand{\twosquares}[1]{\pic{%
  \foreach \i in {1,...,#1} \foreach \j in {1,...,#1} {
    \fill (0.25*\i,0.25*\j) circle (0.06);
    \fill (0.25*\i+0.25*#1+0.3,0.25*\j) circle (0.06);}}}
% A square of n by n dots with the top right dot taken away. A hollow
% circle marks where the dot was.
\newcommand{\squareminusone}[1]{\pic{%
  \foreach \i in {1,...,#1} \foreach \j in {1,...,#1} {
    \ifnum\i=#1 \ifnum\j=#1 \draw[gray] (0.25*\i,0.25*\j) circle (0.06);
      \else \fill (0.25*\i,0.25*\j) circle (0.06);\fi
    \else \fill (0.25*\i,0.25*\j) circle (0.06);\fi}}}
% A square of n by n dots drawn as L-shapes, alternately filled and hollow.
\newcommand{\gnomonsquare}[1]{\pic{%
  \foreach \i in {1,...,#1} \foreach \j in {1,...,#1} {
    \pgfmathtruncatemacro{\m}{max(\i,\j)}
    \ifodd\m \fill (0.25*\i,0.25*\j) circle (0.065);
    \else \draw[thick] (0.25*\i,0.25*\j) circle (0.055);\fi}}}

% A staircase of dots: column k has k dots.
\newcommand{\staircase}[1]{\pic{%
  \foreach \i in {1,...,#1} \foreach \j in {1,...,\i}
    \fill (0.25*\i,0.25*\j) circle (0.06);}}
% Two staircases of n columns, one filled and one hollow, that fit
% together into a rectangle of n by n+1 dots.
\newcommand{\staircasepair}[1]{\pic{%
  \pgfmathtruncatemacro{\top}{#1+1}
  \foreach \i in {1,...,#1} {
    \foreach \j in {1,...,\i} {\fill (0.25*\i,0.25*\j) circle (0.065);}
    \pgfmathtruncatemacro{\s}{\i+1}
    \foreach \j in {\s,...,\top} {\draw[thick] (0.25*\i,0.25*\j) circle (0.055);}}}}

% ----------------------------------------------------------------
% Pixel art. A pixel is a small square, drawn with a white edge so
% that every pixel in a block can be counted. \px[<style>]{col}{row}
% has its top left corner at the column and row given, counting rows
% down from the top.
\tikzset{
  pixel/.style={fill=black!65,draw=white,line width=0.35pt},
  photo/.style={fill=gray!20,draw=white,line width=0.35pt},
  moved/.style={fill=gray!45,draw=white,line width=0.35pt}}
\newcommand{\px}[3][pixel]{\filldraw[#1] ({#2},{-(#3)}) rectangle ++(1,-1);}
% A pale grid of #1 columns by #2 rows for the pixels to sit on.
\newcommand{\canvas}[2]{\draw[gray!35,line width=0.3pt] (0,0) grid (#1,-#2);}
% A picture whose unit is one pixel: \pixpic[<pixel size>]{...}
\newcommand{\pixpic}[2][0.2cm]{\begin{tikzpicture}[x=#1,y=#1,baseline=(current bounding box.south)]#2\end{tikzpicture}}

% A rocket with a body n pixels long.
\newcommand{\rocket}[1]{\pixpic{%
  \pgfmathtruncatemacro{\h}{#1+3}\pgfmathtruncatemacro{\f}{#1+1}\pgfmathtruncatemacro{\g}{#1+2}
  \canvas{5}{\h}
  \px{2}{0}
  \foreach \r in {1,...,#1} \foreach \c in {1,2,3} {\px{\c}{\r}}
  \foreach \c in {0,...,4} {\px{\c}{\f}}
  \px{0}{\g}\px{4}{\g}}}

% A heart of 16 pixels, with every pixel enlarged to an n by n block.
\newcommand{\heart}[1]{\pixpic[0.135cm]{%
  \pgfmathtruncatemacro{\w}{5*#1}
  \canvas{\w}{\w}
  \foreach \c/\r in {1/0,3/0,0/1,1/1,2/1,3/1,4/1,0/2,1/2,2/2,3/2,4/2,1/3,2/3,3/3,2/4}
    \foreach \i in {1,...,#1} \foreach \j in {1,...,#1}
      {\px{\c*#1+\i-1}{\r*#1+\j-1}}}}

% A pyramid of n rows of pixels, with 1, 3, 5, ... pixels in the rows.
% \pyramid[moved]{n} shades the pixels left of the middle column.
\newcommand{\pyramid}[2][pixel]{\pixpic{%
  \pgfmathtruncatemacro{\last}{#2-1}\pgfmathtruncatemacro{\w}{2*#2-1}
  \canvas{\w}{#2}
  \foreach \r in {0,...,\last} {
    \pgfmathtruncatemacro{\a}{\last-\r}\pgfmathtruncatemacro{\b}{\last+\r}
    \foreach \c in {\a,...,\b} {
      \ifnum\c<\last \px[#1]{\c}{\r}\else \px{\c}{\r}\fi}}}}
% The pixels of the pyramid of n rows made into an n by n square: the
% shaded pixels are the ones from the left of the pyramid, turned round.
\newcommand{\pyramidsquare}[1]{\pixpic{%
  \pgfmathtruncatemacro{\last}{#1-1}
  \foreach \c in {0,...,\last} \foreach \r in {0,...,\last} {
    \ifnum\r<\c \px[moved]{\c}{\r}\else \px{\c}{\r}\fi}}}

% A frame one pixel wide around a square picture of n by n pixels.
\newcommand{\pixframe}[1]{\pixpic{%
  \pgfmathtruncatemacro{\last}{#1+1}
  \foreach \c in {0,...,\last} \foreach \r in {0,...,\last} {
    \ifnum\c=0 \px{\c}{\r}\else\ifnum\c=\last \px{\c}{\r}\else
    \ifnum\r=0 \px{\c}{\r}\else\ifnum\r=\last \px{\c}{\r}\else
    \px[photo]{\c}{\r}\fi\fi\fi\fi}}}

% Patterns 1, 2 and 3 of a picture sequence, labelled underneath.
\newcommand{\threepatterns}[3]{%
  \begin{tabular}{@{}c@{\hspace{11pt}}c@{\hspace{11pt}}c@{}}
    #1 & #2 & #3\\[-1pt]
    {\scriptsize Pattern 1} & {\scriptsize Pattern 2} & {\scriptsize Pattern 3}
  \end{tabular}}

% Dotted paper, or squared paper for pixel art, to draw a sequence on.
\newcommand{\dotpaper}[2]{%
  \begin{tikzpicture}[baseline=(current bounding box.north)]
    \foreach \x in {0.25,0.75,...,#1} \foreach \y in {0.25,0.75,...,#2}
      \fill[gray!60] (\x,\y) circle (0.024);
    \draw[gray!60] (0,0) rectangle (#1,#2);
  \end{tikzpicture}}
\newcommand{\pixelpaper}[2]{%
  \begin{tikzpicture}[baseline=(current bounding box.north)]
    \draw[gray!30,line width=0.3pt,step=0.3] (0,0) grid (#1,#2);
    \draw[gray!60] (0,0) rectangle (#1,#2);
  \end{tikzpicture}}

% ----------------------------------------------------------------
% A table of always, sometimes or never statements, carrying on the
% question numbers. \asnrow[<picture>]{<statement>}
\newcounter{asn}
\newcommand{\pill}[1]{\tikz[baseline=(x.base)]\node[draw=gray!70,rounded corners=6pt,
  inner sep=2.5pt,font=\small] (x) {#1};}
\newcommand{\pills}{\pill{Always}\ \pill{Sometimes}\ \pill{Never}}
\newenvironment{asntable}{%
  \setlength{\tabcolsep}{4pt}%
  \begin{tabular}{|>{\raggedright\arraybackslash}m{8.6cm}|m{8.9cm}|}\hline
  \textbf{Statement} & \textbf{Reason}\\\hline}
  {\end{tabular}}
\newcommand{\asnrow}[2][]{\stepcounter{asn}%
  \rule{0pt}{4mm}\textbf{Q\theasn.} #2\par\vspace{3pt}\pills\hfill #1
  & \rule{0pt}{1.4cm}\\\hline}

\begin{document}

\begin{center}
  {\Large\bfseries Picture Sequences and the \boldmath$n$th Term}\\[1pt]
  {\small Extension}
\end{center}
\vspace{-10pt}

\parttitle{1}{Linear sequences}

\begin{keyidea}
A sequence that goes up or down by the same amount each time is a
\textbf{linear sequence}. Its $n$th term is a times table with a number added
or subtracted, such as $3n+1$. We find it by splitting each pattern into the
part that is the same in every pattern and the groups that are added.
\end{keyidea}

\begin{wex}{a row of squares}
\splitlayout{\threepatterns{\chainsquares{1}}{\chainsquares{2}}{\chainsquares{3}}}%
  {Pattern 3, split up:}{\explodedsquares{3}}%
  {\begin{pictable}\filledrow{Matchsticks}{4}{7}{10}{13}{16}{31}{$3n+1$}\end{pictable}}%
  {Each pattern has \textbf{1} matchstick at the start, and $n$ groups of
  \textbf{3} matchsticks. So pattern $n$ has $3\times n+1$ matchsticks, and the
  $n$th term is $3n+1$.}
\end{wex}

\begin{enumerate}
\item Use the same method for this row of houses.\par
\splitlayout{\threepatterns{\chainhouses{1}}{\chainhouses{2}}{\chainhouses{3}}}%
  {Pattern 3, split up:}{\explodedhouses{3}}%
  {\begin{pictable}\picrow{Matchsticks}\end{pictable}}%
  {Each pattern has \gap\ matchstick at the start, and $n$ groups of
  \gap\ matchsticks. So pattern $n$ has \ansline[1.6cm] matchsticks, and the
  $n$th term is \ansline[1.6cm].}

\item Tables in a caf\'e are pushed together in a row. Each dot is a chair.\par
\splitlayout{\threepatterns{\tablesrow{1}}{\tablesrow{2}}{\tablesrow{3}}}%
  {Draw rings round the groups:}{\scalebox{1.5}{\tablesrow{3}}}%
  {\begin{pictable}\picrow{Chairs}\end{pictable}}%
  {Use your rings to explain the $n$th term. \hrulefill}

\item Complete each table and write the $n$th term. Then use the picture to explain the $n$th term.
\begin{enumerate}
\item \pictureandtable{\threepatterns{\dotcross{1}}{\dotcross{2}}{\dotcross{3}}}%
  {\begin{pictable}\picrow{Dots}\end{pictable}}\par
  Explain: \hrulefill
\item \pictureandtable{\threepatterns{\dotL{1}}{\dotL{2}}{\dotL{3}}}%
  {\begin{pictable}\picrow{Dots}\end{pictable}}\par
  Explain where the $-1$ comes from: \hrulefill
\end{enumerate}

\item Use the $n$th terms from the worked example, Q1 and Q2.
\begin{enumerate}
\item Which pattern of the row of squares uses 61 matchsticks?\par
  $3n+1=61$, \ so \ $3n=$ \gap, \ so \ $n=$ \gap
\item How many tables are needed to seat 30 people? \hfill \ansline[3cm]
\item Can a row of houses use exactly 100 matchsticks? Explain. \hfill \ansline[6cm]
\end{enumerate}

\newpage
\item Draw patterns 1, 2 and 3 of a picture sequence for each $n$th term.\par
\begin{tabular}{@{}l@{\hspace{8pt}}l@{\hspace{8pt}}l@{}}
(a) \ $2n+3$ & (b) \ $4n-1$ & (c) \ $5n+2$\\
\dotpaper{5.5}{2.3} & \dotpaper{5.5}{2.3} & \dotpaper{5.5}{2.3}
\end{tabular}
\end{enumerate}

\parttitle{2}{Square numbers}

\begin{keyidea}
$n^2$, read as ``$n$ squared'', means $n\times n$. It is the number of dots in
a square with $n$ dots along each side. The numbers $1, 4, 9, 16, \ldots$ are the
\textbf{square numbers}.
\end{keyidea}

\begin{enumerate}[resume]
\item
\begin{enumerate}
\item Complete the table for the squares of dots.\par
\pictureandtable{\threepatterns{\dotsquare{1}}{\dotsquare{2}}{\dotsquare{3}}}%
  {\begin{pictable}\picrow{Dots}\end{pictable}}
\item \begin{minipage}[t]{0.8\linewidth}
  Pattern 5 is drawn as L-shapes of filled and hollow dots. Write the number of
  dots in each L-shape.\par\smallskip
  $25 = 1 + {}$\gap${}+{}$\gap${}+{}$\gap${}+{}$\gap\par\smallskip
  The L-shape that is added to make pattern $n$ has \ansline[2cm] dots.
  \end{minipage}\hfill
  \begin{minipage}[t]{0.15\linewidth}\raggedleft\vspace{0pt}\gnomonsquare{5}\end{minipage}
\end{enumerate}

\item Each pattern is a square of dots with dots added or taken away. Count the
filled dots, complete the table and write the $n$th term.
\begin{enumerate}
\item \pictureandtable{\threepatterns{\squareplustwo{1}}{\squareplustwo{2}}{\squareplustwo{3}}}%
  {\begin{pictable}\picrow{Dots}\comparerow{$n^2$}\end{pictable}}
\item \pictureandtable{\threepatterns{\dotrect{1}{2}}{\dotrect{2}{3}}{\dotrect{3}{4}}}%
  {\begin{pictable}\picrow{Dots}\comparerow{$n^2$}\end{pictable}}
\item \pictureandtable{\threepatterns{\twosquares{1}}{\twosquares{2}}{\twosquares{3}}}%
  {\begin{pictable}\picrow{Dots}\end{pictable}}
\item \pictureandtable{\threepatterns{\squareminusone{1}}{\squareminusone{2}}{\squareminusone{3}}}%
  {\begin{pictable}\picrow{Dots}\end{pictable}}
\end{enumerate}

\item These dots make staircases.
\begin{enumerate}
\item Complete the table. Leave the $n$th term until part (c).\par
\pictureandtable{\threepatterns{\staircase{1}}{\staircase{2}}{\staircase{3}}}%
  {\begin{pictable}\picrow{Dots}\end{pictable}}
\item \begin{minipage}[t]{0.8\linewidth}
  Two copies of pattern 4 fit together to make a rectangle of 4 by 5 dots.
  How many dots are in the rectangle made from two copies of pattern $n$?
  \ansline[2.6cm]
  \end{minipage}\hfill
  \begin{minipage}[t]{0.15\linewidth}\raggedleft\vspace{0pt}\staircasepair{4}\end{minipage}
\item Use your answer to (b) to write the $n$th term in the table.
\end{enumerate}

\item Draw patterns 1, 2 and 3 of a picture sequence for each $n$th term.\par
\begin{tabular}{@{}l@{\hspace{8pt}}l@{\hspace{8pt}}l@{}}
(a) \ $n^2+3$ & (b) \ $3n^2$ & (c) \ $n^2+2n$\\
\dotpaper{5.5}{2.3} & \dotpaper{5.5}{2.3} & \dotpaper{5.5}{2.3}
\end{tabular}
\end{enumerate}

\newpage
\parttitle{3}{Pixel art sequences}

\begin{enumerate}[resume]
\item A pixel rocket gets longer in each pattern.
\begin{enumerate}
\item Count the pixels, complete the table and write the $n$th term.\par
\pictureandtable{\threepatterns{\rocket{1}}{\rocket{2}}{\rocket{3}}}%
  {\begin{pictable}\picrow{Pixels}\end{pictable}}
\item Which rocket has 50 pixels? \ansline[2cm]
\item Can a rocket have exactly 100 pixels? Explain. \hrulefill
\end{enumerate}

\item Pattern 1 is a pixel heart. In pattern 2 every pixel of the heart becomes a
2 by 2 block, and in pattern 3 every pixel becomes a 3 by 3 block.
\begin{enumerate}
\item Complete the table and write the $n$th term.\par
\pictureandtable{\threepatterns{\heart{1}}{\heart{2}}{\heart{3}}}%
  {\begin{pictable}\picrow{Pixels}\end{pictable}}
\item Pattern 2 is twice as wide as pattern 1. Explain why it has 4 times as
many pixels, and not 2 times as many. \hrulefill
\end{enumerate}

\item Each pattern is a pyramid of pixels.
\begin{enumerate}
\item Complete the table and write the $n$th term.\par
\pictureandtable{\threepatterns{\pyramid{1}}{\pyramid{2}}{\pyramid{3}}}%
  {\begin{pictable}\picrow{Pixels}\end{pictable}}
\item \begin{minipage}[t]{0.62\linewidth}
  In pattern 4 the shaded pixels are moved to turn the pyramid into a square.
  Use the picture to explain why the pyramid in pattern $n$ has $n^2$ pixels.\par
  \vspace{10pt}\hrulefill\par\vspace{10pt}\hrulefill
  \end{minipage}\hfill
  \begin{minipage}[t]{0.34\linewidth}\raggedleft\vspace{0pt}%
  \pyramid[moved]{4}\ \raisebox{0.35cm}{$\longrightarrow$}\ \pyramidsquare{4}\end{minipage}
\end{enumerate}

\item A pixel artist draws a frame one pixel wide around a square picture of
$n$ by $n$ pixels.
\begin{enumerate}
\item Complete the table and write each $n$th term.\par
\pictureandtable{\threepatterns{\pixframe{1}}{\pixframe{2}}{\pixframe{3}}}%
  {\begin{pictable}\picrow{Picture}\picrow{Frame}\picrow{Total}\end{pictable}}
\item Use the picture to explain the $n$th term of the frame. \hrulefill
\end{enumerate}

\item \begin{minipage}[t]{0.44\linewidth}
  Design a pixel art sequence of your own, such as a snake, a tree or a robot
  that grows. Draw patterns 1, 2 and 3, then find the $n$th term. For a
  challenge, design one with the $n$th term $n^2+4$.
  \end{minipage}\hfill
  \begin{minipage}[t]{0.55\linewidth}\raggedleft\vspace{0pt}
  \begin{pictable}\picrow{Pixels}\end{pictable}\end{minipage}\par\smallskip
\pixelpaper{17.1}{3.6}
\setcounter{asn}{\value{enumi}}
\end{enumerate}

\newpage
\parttitle{4}{Always, sometimes or never?}

Decide whether each statement is always true, sometimes true or never true,
and circle your answer. Explain how you know in the reason box. For sometimes,
give one example where the statement is true and one where it is false. In
every statement $n$ is a pattern number, so $n$ is 1, 2, 3, \ldots

\begin{asntable}
\asnrow{A term of $3n+1$ is one more than a multiple of 3.}
\asnrow{A term of $4n+2$ is a multiple of 4.}
\asnrow{A term of $2n+5$ is a multiple of 5.}
\asnrow{If an $n$th term ends in $+1$, like $2n+1$ or $3n+1$, then every term is odd.}
\asnrow{$n^2$ is greater than $2n$.}
\asnrow[\dotrect{3}{4}]{$n^2+n$ is even.}
\asnrow{A square number ends in the digit 7.}
\asnrow{$n^2+1$ is a square number.}
\asnrow[\gnomonsquare{4}]{Two square numbers next to each other have an odd difference.}
\asnrow[\heart{1}\quad\heart{2}]{Doubling the pattern number of the pixel heart doubles the number of pixels.}
\asnrow[\pixframe{1}\quad\pixframe{2}]{The number of pixels in the frame around a square picture is a multiple of 4.}
\asnrow[\staircase{1}\quad\staircase{2}\quad\staircase{3}]{The number of dots in a staircase is odd.}
\end{asntable}

\vspace{4pt}
\stepcounter{asn}\textbf{Q\theasn.} Write your own statement of each type about $n$th terms.\\[7pt]
Always: \hrulefill\\[7pt]
Sometimes: \hrulefill\\[7pt]
Never: \hrulefill

\end{document}
